\begin{corolario}[Coste de salida en el tramo mesoscópico]\label{coste-debil}
Para un tablero de primera salida del tramo mesoscópico, que cumple las cotas de la clase exterior y
en el que $\max_f\ell_f\le\varepsilon_{\rm ext}<1$, se tiene, para $k$ suficientemente grande,
\[
 g^+(B)\le C_0+\beta(\varepsilon_{\rm ext})k,\qquad
 \beta(v)=\frac{2v}{1-v},
\]
con la función $\beta$ de la Sección~\ref{meso:constantes}. La constante $C_0$ se puede elegir
uniformemente para esas cotas.
\end{corolario}
\begin{proof}
Para $0\le x<1$ se tiene $-\log(1-x)-x=\int_0^x t/(1-t)\,dt\le x^2/(2(1-x))$. En el factor de $F$ que
es común a todas las permutaciones obtenemos, usando $\log(1-\ell_f)\le-\ell_f$,
\[
 \sum_f\log h_f-2\sum_e\log(1-p_e)
 \le-\sum_f\ell_f+2\sum_ep_e+
       \frac{\sum_ep_e^2}{1-p_{\max}}
 =O_C(1).
\]
Para las permutaciones legales el numerador $(1-p_e)^2$ de cada factor de casilla es a lo sumo uno.
Además $-\log h_f\le\ell_f/(1-\ell_f)\le\varepsilon_{\rm ext}/(1-\varepsilon_{\rm ext})$. Hay dos de
esos logaritmos por cada una de las $k$ casillas elegidas. Se sigue que
$F(\sigma)\le\exp(C'+\beta(\varepsilon_{\rm ext})k)$, desigualdad que también vale si el indicador de
legalidad es cero. Las cotas exteriores dan \eqref{terminal-hipotesis-matriz} con $K=k$,
$M_0=C_{\rm ext}$ y $\rho_*=1-\gamma_{\rm ext}$, por el párrafo que sigue a
\eqref{terminal-hipotesis-matriz}. El Lema~\ref{coste-factor} y la
Proposición~\ref{coste-normalizadores} dan la cota
para $g$ cuando $Q(B)>0$. Se absorben la cota uniforme de $\log(\Gamma_4/\Gamma_2)$ y el $o(1)$
uniforme en $C_0$, que se toma positivo. Esto permite pasar a $g^+$, y cubre asimismo el caso
$Q(B)=0$.
\end{proof}
